%! Sine and Cosine Function Values — Unit 3, Lesson 3.3 — Homework (Answer Key)
\documentclass[10pt]{article}
\usepackage{precalculus-article}
\usepackage{precalculus-key}   % pulls in -boxes; do NOT also load -boxes
\newcommand{\TallMath}[1]{$\displaystyle #1\rule[-1.4em]{0pt}{3.2em}$}

\begin{document}

\pageheader{Unit 3, Lesson 3.3}{Homework --- Answer Key}
\namedateperiod

\vspace{0.1in}

\begin{notesbox}{Sine and Cosine Function Values --- Practice}

\textbf{1.} Complete the table with exact values.

\vspace{0.06in}
\renewcommand{\arraystretch}{2.2}
\begin{tabularx}{\linewidth}{@{} >{\bfseries}l *{6}{>{\centering\arraybackslash}X} @{}}
\toprule
$\theta$ (rad) & $\dfrac{5\pi}{4}$ & $\dfrac{4\pi}{3}$ & $\dfrac{3\pi}{2}$ & $\dfrac{5\pi}{3}$ & $\dfrac{7\pi}{4}$ & $\dfrac{11\pi}{6}$ \\
\midrule
$\sin\theta$ &
  \ans{$-\dfrac{\sqrt{2}}{2}$} &
  \ans{$-\dfrac{\sqrt{3}}{2}$} &
  \ans{$-1$} &
  \ans{$-\dfrac{\sqrt{3}}{2}$} &
  \ans{$-\dfrac{\sqrt{2}}{2}$} &
  \ans{$-\dfrac{1}{2}$} \\
$\cos\theta$ &
  \ans{$-\dfrac{\sqrt{2}}{2}$} &
  \ans{$-\dfrac{1}{2}$} &
  \ans{$0$} &
  \ans{$\dfrac{1}{2}$} &
  \ans{$\dfrac{\sqrt{2}}{2}$} &
  \ans{$\dfrac{\sqrt{3}}{2}$} \\
\bottomrule
\end{tabularx}
\renewcommand{\arraystretch}{1}

\vspace{0.12in}

\textbf{2.} $r = 8$, $\theta = 7\pi/6$. Reference angle $= \pi/6$ (QIII: both negative).
\[
  x = 8\cos\!\left(\tfrac{7\pi}{6}\right) = 8\cdot\left(-\tfrac{\sqrt{3}}{2}\right) = -4\sqrt{3},
  \qquad
  y = 8\sin\!\left(\tfrac{7\pi}{6}\right) = 8\cdot\left(-\tfrac{1}{2}\right) = -4
\]

$P = \Bigl($ \ans{$-4\sqrt{3}$} $,$ \ans{$-4$} $\Bigr)$

\vspace{0.12in}

\textbf{3.} $\sin\theta = -\sqrt{2}/2$ and $\cos\theta < 0$: negative sine and negative cosine $\Rightarrow$ QIII.
Reference angle $= \pi/4$, so $\theta = \pi + \pi/4$.

$\theta = $ \ans{$\dfrac{5\pi}{4}$}

\vspace{0.12in}

\textbf{4.} $\theta = 5\pi/3$ is in QIV. Reference angle $= \pi/3$.
$\sin(\pi/3) = \sqrt{3}/2$; negate for QIV: $\sin(5\pi/3) = -\sqrt{3}/2$.
$y$-coordinate $= r\sin(5\pi/3) = -r\cdot\sqrt{3}/2$.

Answer: \ans{(B)} \quad Justify: \ans{$5\pi/3$ is in QIV (ref.\ angle $\pi/3$); $\sin$ is negative in QIV, so $y = -r\cdot\sqrt{3}/2$.}

\vspace{0.12in}

\textbf{5.} On the unit circle, the point at angle $\theta$ is $(x, y) = (\cos\theta, \sin\theta)$.
Adding $\pi$ to an angle rotates the point $180^\circ$, reflecting it through the origin.
This maps $(x, y)$ to $(-x, -y)$.
Therefore the $x$-coordinate at $\pi + \theta$ is $-\cos\theta$.

\ans{Because rotating by $\pi$ reflects the terminal point through the origin, the $x$-coordinate changes sign: $\cos(\pi+\theta) = -\cos\theta$.}

\end{notesbox}

\begin{extensionbox}
\textbf{Extension:} $\sin(\pi - \theta) = \sin\theta$.

\ans{The point at angle $\theta$ is $(\cos\theta, \sin\theta)$. Replacing $\theta$ with $\pi - \theta$ reflects the point across the $y$-axis (the $x$-coordinate negates but the $y$-coordinate is unchanged). Therefore $\sin(\pi - \theta) = \sin\theta$.}
\end{extensionbox}

\vspace{0.1in}

\begin{tcolorbox}[colback=goldbg, colframe=goldacc, arc=2mm,
    left=3mm, right=3mm, top=2mm, bottom=2mm]
\small\textbf{Preview of Lesson 3.4:} Now that we know exact $\sin\theta$ and $\cos\theta$ values
at key angles, we will use those values as ordered pairs to build and interpret the
graphs of $y = \sin x$ and $y = \cos x$. Bring your completed homework table --- those
values become the points on the graph!
\end{tcolorbox}

\end{document}
